Region $R_2$ is dealt with by choosing $\delta_2$ sufficiently small that
$I_2(\phi)$ is dominated by a particular Gaussian integral.  The we will make
that Gaussian integral small by choosing $\delta_1$ sufficiently large.

\begin{lem}\lemlabel{region_r2}
%
Let \assuref{core_bclt_assu} hold. Take $\delta_1 \ge \sqrt{8 / \infoev}$.
There exist $\delta_2$ and $\epsilon_U$, defined as functions only of the
population quantities $\infoev$ and $\likk{k}(\theta)$, such that, for any
$\xvec$ such that $\regevent$ occurs,
%
\begin{align*}
%
\norm{I_2(\phi)}_2 \le&
    O\left(N^{-3}\right)
    \expect{\normresid{\infoev / 2}{R_1 \bigcup R_2}(\tau)}
           {\norm{\phi(\theta, \xvec_s)}_2 }.
\end{align*}
%
\end{lem}
%
\begin{proof} %%%%%%%%%%%%% proof of region_r2 lemma
%
The main part of the proof consists in showing that there exists a $\delta_2$
sufficiently small so that
%
\begin{align}\eqlabel{r2_target_bound}
%
\norm{I_2(\phi)}_2 \le& \int_{R_2}
    \exp\left(
        -\frac{\infoev}{2} \norm{\tau}_2^2
        \right) \norm{\phi(\theta, \xvec_s)}_2 d\tau,
%
\end{align}
%
The approach will be to bound the difference between $N \likhat(\theta) - N
\likhat(\thetahat)$ and $\frac{1}{2} \likk{2}(\thetatrue)$ so that the entire
expression inside the integral's exponent can be written as  $A \tau^2$ for a
positive definite matrix $A$.

For the duration of this proof, define the following residual matrix:
%
\begin{align*}
\Lambda(\theta, \xvec, N) :={}&
    \frac{1}{2} \left(\likhatk{2}(\thetahat) - \likk{2}(\thetatrue)\right) + \\
&   N^{-1/2} \frac{1}{6}\left(
            \likhatk{3}(\thetahat) -
            \likk{3}(\thetatrue)
        \right) \tau + \\
&    N^{-1/2} \frac{1}{6}\likk{3}(\thetatrue) \tau +
    N^{-1} \tresid{4}{\likhat, \thetahat, \theta} \tau^2.
\end{align*}
%
Under this definition, by \eqref{taylor_expansion},
%
\begin{align*}
%
N \likhat(\theta) - N \likhat(\thetahat) ={}&
\left( \frac{1}{2} \likk{2}(\thetatrue) + \Lambda(\theta, \xvec, N) \right) \tau^2.
%
\end{align*}
%
This expansion holds for all $\theta$, but we need to control the
residual only in $R_2$.  We do so by bounding the entries of the
$\thetadim \times \thetadim$ matrix $\Lambda(\theta, \xvec,
N)$.  When $\regevent$ occurs, we have
%
\begin{align*}
%
\norm{\likhatk{2}(\thetahat) - \likk{2}(\thetatrue)}_2 \le&
\norm{\likhatk{2}(\thetahat) - \likk{2}(\thetahat)}_2 +
\norm{\likk{2}(\thetahat) - \likk{2}(\thetatrue)}_2 \\
\le&
\sup_{\theta \in R_2}
    \norm{\likhatk{2}(\theta) - \likk{2}(\theta)}_2 +
\sup_{\theta \in R_2} \norm{\likk{3}(\theta)}_2
    \norm{\thetahat - \thetatrue}_2 \\
\le& \left(1 + \sup_{\theta \in R_2} \norm{\likk{3}(\theta)}_2\right)
    \epsilon_U.
\quad\textrm{(\lemref{regular_event}, \eventref{ulln, thetahat_close})}
%
\end{align*}
%
Similarly,
%
\begin{align*}
%
\norm{\left(\likhatk{3}(\thetahat) - \likk{3}(\thetatrue)\right) \tau}_2
\le&
\norm{\likhatk{3}(\thetahat) - \likk{3}(\thetatrue)}_2
\norm{\tau}_2
\\
\le& \left(1 + \sup_{\theta \in R_2} \norm{\likk{4}(\theta)}_2\right)
    \norm{\tau}_2 \epsilon_U.
%
\end{align*}
%
Next, applying Cauchy-Schwarz gives
%
\begin{align*}
%
\norm{\likk{3}(\thetatrue) \tau}_2 \le
    \norm{\likk{3}(\thetatrue)}_2 \norm{\tau}_2.
%
\end{align*}
%
Finally, let us consider the Taylor series residual term.  Let $\thetatil(t) = t
(\theta - \thetahat) + \thetahat$.  Because we can write
%
\begin{align*}
\tresid{4}{\likhat, \thetahat, \theta} ={}&
\frac{1}{6} \int_0^1 (t - 1)^3 \likhatk{4}(\thetatil(t)) dt \\
={}&
\frac{1}{6} \int_0^1 (t - 1)^3 \left(
    \likhatk{4}(\thetatil(t)) - \likk{4}(\thetatil(t))
\right) dt +
\frac{1}{6} \int_0^1 (t - 1)^3 \likk{4}(\thetatil(t)) dt,
\end{align*}
%
we have, when $\regevent$ occurs,
%
\begin{align*}
\sup_{\theta \in R_2}
\norm{\tresid{4}{\likhat, \thetahat, \theta} \tau^2 }_2 \le&
\norm{\tresid{4}{\likhat, \thetahat, \theta}}_2 \norm{\tau}_2^2 \\
\le&
\frac{1}{6} \left(
    \sup_{\theta \in R_2}
        \norm{\likhatk{4}(\theta) - \likk{4}(\theta)}_2 +
    \sup_{\theta \in R_2} \norm{\likk{4}(\theta)}_2 \right) \norm{\tau}_2^2 \\
\le&
\frac{1}{6} \left(
    \epsilon_U + \sup_{\theta \in R_2} \norm{\likk{4}(\theta)}_2 \right)
    \norm{\tau}_2^2.
\quad\textrm{(\lemref{regular_event}, \eventref{ulln})}
\end{align*}
%

Combining the above results, together with the fact that, in $R_2$, $\sqrt{N}
\norm{\tau}_2 \le \delta_2$, gives
%
\begin{align}
%
\sup_{\theta \in R_2} \norm{\Lambda(\theta, \xvec, N)}_2 \le&
    \frac{1}{2} \norm{\likhatk{2}(\thetahat) - \likk{2}(\thetatrue)}_2 +
        \nonumber\\&
    N^{-1/2} \frac{1}{6}\norm{
        \likhatk{3}(\thetahat) -
        \likk{3}(\thetatrue)
    } \sup_{\theta \in R_2} \norm{\tau}_2 +
        \nonumber\\&
    N^{-1/2} \frac{1}{6} \norm{\likk{3}(\thetatrue)}_2
        \sup_{\theta \in R_2} \norm{\tau}_2 +
        \nonumber\\&
    N^{-1} \sup_{\theta \in R_2}
        \norm{\tresid{4}{\likhat, \thetahat, \theta}}_2 \norm{\tau}_2^2
\nonumber\\\le&
    \frac{1}{2} \left(1 + \sup_{\theta \in R_2} \norm{\likk{3}(\theta)}_2\right)
        \epsilon_U +
   \frac{1}{6} \left(1 + \sup_{\theta \in R_2} \norm{\likk{4}(\theta)}_2\right)
    \delta_2 \epsilon_U + \nonumber\\
&   \frac{1}{6} \norm{\likk{3}(\thetatrue)}_2 \delta_2 +
   \frac{1}{6} \left(
        \epsilon_U + \sup_{\theta \in R_2} \norm{\likk{4}(\theta)}_2 \right)
    \delta_2^2. \eqlabel{r2_lambda_bound}
%
\end{align}
%
The preceding bound contains only the $\epsilon_U$, $\delta_2$, and the
population quantities $\sup_{\theta \in R_2} \norm{\likk{3}(\theta)}_2$,
$\sup_{\theta \in R_2} \norm{\likk{4}(\theta)}_2$, and $\infoev$. The quantities
$\sup_{\theta \in R_2} \norm{\likk{3}(\theta)}_2$ and $\sup_{\theta \in R_2}
\norm{\likk{4}(\theta)}_2$ depend on $\delta_2$ themselves, but are finite and
decreasing in $\delta_2$.

Consequently, when $\regevent$ occurs, there exist an $\epsilon_U$ and
$\delta_2$ sufficiently small so that
%
\begin{align}\eqlabel{re2_lambda_inf_bound}
%
\norm{\Lambda(\theta, \xvec, N)}_\infty \le
    \sqrt{\thetadim^2} \norm{\Lambda(\theta, \xvec, N)}_2 \le \frac{\infoev}{2},
%
\end{align}
%
independently of $\xvec$ (except that $\regevent$ occurs). The bound
\eqref{re2_lambda_inf_bound} cannot necessarily be achieved simply by
setting $\epsilon_U$ to be small (i.e., by relying on the consistency of
$\thetahat$ and the ULLN) because of the terms $\sup_{\theta \in R_2}
\norm{\likk{4}(\theta)}_2 \delta_2^2$ and $\norm{\likk{3}(\thetatrue)}_2
\delta_2$ in \eqref{r2_lambda_bound}---it is necessary to set $\delta_2$
small as well.
% However, using the bound
% \eqref{r2_lambda_bound}, when can choose $\epsilon_U$ and $\delta_2$ as a
% function only of population quantities to guarantee \eqref{re2_lambda_inf_bound}
% whenever $\regevent$ occurs.

Recall from \defref{loglik_def} that $\info =
-\frac{1}{2}\likk{2}(\thetatrue)$, and from \defref{loglik_details_def} that
$\infoev > 0$ is the minimum eigenvalue of $\info$.  These facts, combined
with \eqref{re2_lambda_inf_bound} give that, for any $\tau$,
%
\begin{align*}
%
\sup_{\theta \in R_2}
\left(\frac{1}{2}\likk{2}(\thetatrue) + \Lambda(\theta, \xvec, N)\right) \tau^2
    \le& -\infoev \norm{\tau}_2^2 + \frac{\infoev}{2} \norm{\tau}_2^2 \\
    ={}& -\frac{\infoev}{2} \norm{\tau}_2^2.
%
\end{align*}
%
Plugging into the definition of $I_2(\phi)$ given in \defref{region_division}
gives the desired \eqref{r2_target_bound}.

To complete the proof, we use the fact that the normal density is bounded
in $R_2$.

\begin{align*}
%
\norm{I_2(\phi)}_2
\le&
\int_{R_2}
    \exp\left(
        -\frac{\infoev}{2} \norm{\tau}_2^2
        \right) \norm{\phi(\theta, \xvec_s)}_2  d\tau
\\={}&
\int_{R_2}
    \exp\left(
        -\frac{\infoev}{4} \norm{\tau}_2^2
        -\frac{\infoev}{4} \norm{\tau}_2^2
        \right) \norm{\phi(\theta, \xvec_s)}_2 d\tau
\\\le&
\sup_{\tau \in R_2}
\exp\left(
    -\frac{\infoev}{4} \norm{\tau}_2^2
\right)
\int_{R_2}
    \exp\left(
        -\frac{\infoev}{4} \norm{\tau}_2^2
        \right) \norm{\phi(\theta, \xvec_s)}_2 d\tau
\\\le &
\exp\left(
    -\frac{\infoev}{4} \delta_1^2 \log N
\right)
\int_{R_1 \bigcup R_2}
    \exp\left(
        -\frac{\infoev}{4} \norm{\tau}_2^2
        \right) \norm{\phi(\theta, \xvec_s)}_2 d\tau.
%
\end{align*}
%
In the final line, we have used the fact that $(\log N)^2 \ge \log N$ for all $N
\ge 3$. Taking $\delta_1^2 \ge 12 / \infoev$ and multiplying by the $O(1)$
normalizer for the $\normresid{\infoev/2}{R_1 \bigcup R_2}(\tau)$ completes the proof.
%
\end{proof} %%%%%%%%%%%%% proof of region_r2 lemma

By inspecting the proof of \lemref{region_r2}, one can see that that any
polynomial rate of convergence can be achieved by taking $\delta_1$ sufficiently
large. Indeed, if the lower upper boundary of $R_1$ were of the form $\tau \le
N^k$ for some $k < 1/2$, then $\norm{I_2(\phi)}_2$ would decay exponentially and
$R_1$ would still shrink to a point in the $\theta$ domain.  However, such a
rapidly growing boundary between $R_1$ and $R_2$ comes at a cost in the analysis
of region $R_1$.  In particular, it is the use of a logarithmically growing
boundary is what allows us to give polylog bounds in the subsequent
\lemref{region_r1_final_bound} and in the final result of
\thmref{bayes_clt_main}.
